%%%%%%%%%%%%%%%%%%%%%%%%%%%%%%%%%%%%%%%
%% Sustentación Tesis Doctoral        
%% Update: 21 de enero, 2023
%%%%%%%%%%%%%%%%%%%%%%%%%%%%%%%%%%%%%%%

%%%%%%% Initial Setup %%%%%%%
\documentclass[10pt, xcolor=dvipsnames]{beamer}
\usetheme{Boadilla}
\input{preamble/slides}

%%%%%%% Title, Author and Institution %%%%%%%
\title{\textbf{Semana 1: Bienvenida e Introducción}}
\subtitle{\LARGE{Analitica para los negocios} \\ \vspace{0.2cm}
          \large{06278 - ECO \,|\, Periodo 202610}}
\author[]{PhD, Eduard Fernando Martínez González \\ \vspace{0.2cm}
          \normalsize{Universidad Icesi}}
\date{2 de febrero de 2026}

%%%%=================================%%%%
%% Portada
\begin{document} 
\begin{frame}
\titlepage
\end{frame}

%%%%=================================%%%%
\section{Bienvenida y presentaciones}
%%%%=================================%%%%

%%%%=================================%%%%
%% Presentación del profesor
\begin{frame}{Acerca de mí}
\begin{itemize}
  \item \textbf{Eduard F. Martínez-González} (PhD en Economía)
  \item Profesor del Departamento de Economía, Universidad Icesi
  \item Director del CIENFI (Centro de Investigación en Economía y Finanzas).
  \item Intereses: Econometría aplicada,  Big Data y Machine Learning para política pública y toma de decisiones.
  \item Enfoque: investigación aplicada + consultoría
\end{itemize}
%%
\vspace{0.2cm}
\textbf{Horario de atención:} Lunes de 8:30 a 9:30 y Martes de 8:00 a 9:00.
%%
\vspace{0.2cm}
\textbf{Oficina:} Dirección CIENFI (Bloque B)
%%
\vspace{0.2cm}
\textbf{Enlaces y contacto} \\ \vspace{0.1cm}
\href{https://eduard-martinez.github.io/}{Sitio web} \,|\, 
\href{https://github.com/eduard-martinez}{GitHub} \,|\, 
\href{https://x.com/emartigo}{X} \,|\, 
\href{https://www.linkedin.com/in/eduard-fernando-mart\%C3\%ADnez-gonz\%C3\%A1lez-b99937117/}{linkedin} \,|\, 
\href{mailto:efmartinez@icesi.edu.co}{efmartinez@icesi.edu.co}
\end{frame}


%%%%=================================%%%%
%% Monitores
\begin{frame}{Equipo de apoyo (Monitorías y contacto)}
%%
\begin{itemize}
  \item \textbf{Miguel Angel Torres Montoya}
  \begin{itemize}
    \item Correo: \href{mailto:miguel.torres2@u.icesi.edu.co}{miguel.torres2@u.icesi.edu.co}
    \item Tel/WhatsApp: +57 316 244 0219
  \end{itemize} \vspace{0.25cm}
%%
  \item \textbf{Maria Paula Ceballos}
  \begin{itemize}
    \item Correo: \href{mailto:paulaceballostol@gmail.com}{paulaceballostol@gmail.com}
    \item Tel/WhatsApp: +57 301 217 5806
  \end{itemize}
\end{itemize} \vspace{0.25cm}
%%
\footnotesize
\textit{Uso recomendado: dudas de talleres, soporte con R, revisión de errores, y acompañamiento para el proyecto final.}
\end{frame}

%%%%=================================%%%%
%% Conozcámonos (30--45s por persona)
\begin{frame}{Conozcámonos (30--45s por persona)}
\textbf{\large Preguntas guía:}  \vspace{0.2cm}
\begin{itemize}
  \item Nombre y programa/semestre. \vspace{0.2cm}
    \vspace{0.2cm}
  \item En una escala 1--5: ¿cómo te sientes con programación?  \vspace{0.2cm}
  \item ¿Has usado R u otra herramienta de datos? (Excel, Python, Stata, Power BI, etc.)  \vspace{0.2cm}
  \item ¿Qué te gustaría poder hacer con datos al final del curso?  \vspace{0.2cm}
\end{itemize}
\end{frame}

%%%%=================================%%%%
%% Agenda
\begin{frame}[noframenumbering]{Agenda:}
\tableofcontents[]
\end{frame}

%%%%=================================%%%%
\section{¿Por qué Business Analytics?}
\begin{frame}[noframenumbering]{Agenda:}
\tableofcontents[currentsection]
\end{frame}


%%%%=================================%%%%
%% ¿Qué es Business Analytics?
\begin{frame}{¿Qué es Business Analytics?}
\footnotesize
Business Analytics es el uso de \textit{datos + métodos analíticos} para \textit{entender} una situación, \textit{predecir} lo que puede pasar y \textit{recomendar acciones} para tomar mejores decisiones.
\begin{itemize}
  \item \textbf{¿Qué produce?}
  \begin{itemize}
    \item \textit{Insights} (qué está pasando y por qué) + \textit{métricas/KPIs} + \textit{visualizaciones}
    \item \textit{Modelos} (predicción o clasificación) y \textit{experimentos} (A/B testing)
    \item \textit{Recomendaciones} accionables (qué haríamos diferente con evidencia)
  \end{itemize}
  \item \textbf{¿Cómo se hace?}
  \begin{itemize}
    \item Pregunta $\rightarrow$ datos $\rightarrow$ limpieza/EDA $\rightarrow$ análisis/modelo $\rightarrow$ comunicación $\rightarrow$ decisión \,(\textit{y se repite})
  \end{itemize}
\end{itemize} \vspace{-0.4cm}

\begin{figure}
    \centering
    \includegraphics[width=0.5\linewidth]{slides/week-01/figures/que_ba.png}
    \caption{Caption}
    \label{fig:placeholder}
\end{figure}
\end{frame}


%%%%=================================%%%%
%% Slide 1 — Quiero empleo
\begin{frame}{Motivación 1: “Quiero empleo” (mercado laboral)}
\small
\textit{Señal del mercado:} los roles ligados a datos/IA aparecen entre los de mayor crecimiento esperado. \vspace{-0.2cm}
 \begin{figure}
     \centering
     \caption*{Fastest-growing and fastest-declining jobs, 2025-2030}
     \includegraphics[width=0.62\linewidth]{slides/week-01/figures/job_created.png} \\
     \footnotesize \textit{Fuente: \href{https://www.weforum.org/publications/the-future-of-jobs-report-2025/}{World Economic Forum. (2025). \textit{The Future of Jobs Report 2025}}}
     \label{fig:placeholder}
 \end{figure} \vspace{0.1cm}

\textit{Traducción a habilidades:} más allá del cargo, lo que se demanda es \textit{fluidez con datos}: \vspace{-0.2cm}
\begin{itemize}
\item entender un dataset, resumirlo, graficarlo y comunicar hallazgos;
\item construir un pipeline básico: datos $\rightarrow$ modelo $\rightarrow$ validación $\rightarrow$ decisión.
\end{itemize}
\end{frame}


%%%%=================================%%%%
%% Slide 2 — Quiero emprender
\begin{frame}{Motivación 2: “Quiero emprender”}
\small
\textit{Emprender es un problema de decisión bajo incertidumbre (decidir rápido y fallar barato):} Muchísimas startups fracasan por no entender el mercado o por quedarse sin caja.
 \begin{figure}
     \centering
     \includegraphics[width=0.65\linewidth]{slides/week-01/figures/startup_fail_reasons.png} \\
     \footnotesize \textit{Fuente: \href{https://www.cbinsights.com/research/report/startup-failure-reasons-top/}{CB Insights (2021). \textit{The Top 12 Reasons Startups Fail}}}
     \label{fig:placeholder}
 \end{figure} \vspace{0.1cm}

\textit{¿Qué aporta Business Analytics?} un enfoque sistemático para: \vspace{-0.2cm}
  \begin{itemize}
    \item medir adquisición/retención/conversión (qué funciona y qué no).
    \item experimentar (A/B testing) para aprender con evidencia.
  \end{itemize}
\end{frame}

%%%%=================================%%%%
%% Slide 3 — No me dedicaré a datos, pero necesito entender
\begin{frame}{Motivación 3: Alfabetización analítica}
\small
\textit{Aunque no seas “data scientist”,} vas a trabajar con gente, reportes y decisiones basadas en datos. \vspace{-0.2cm}
 \begin{figure}
     \centering
     \caption*{Skills on the rise, 2025-2030}
     \includegraphics[width=0.5\linewidth]{slides/week-01/figures/hh.png} \\
     \footnotesize \textit{Fuente: \href{https://www.weforum.org/publications/the-future-of-jobs-report-2025/}{World Economic Forum. (2025). \textit{The Future of Jobs Report 2025}}}
     \label{fig:placeholder}
 \end{figure} \vspace{0.1cm}
\footnotesize
\textbf{La ventaja competitiva mínima:} entender la jerga y hacer buenas preguntas: \vspace{-0.2cm}
  \begin{itemize}
    \item cuándo un gráfico/tabla “prueba” algo y cuándo no.
    \item entrenamiento vs test, validación cruzada, métricas (accuracy, RMSE, etc.).
  \end{itemize}
\end{frame}


%%%%=================================%%%%
%% Agenda
\section{Acerca de este curso}
\subsection{Contenido}
\begin{frame}[noframenumbering]{Agenda:}
\tableofcontents[currentsection,currentsubsection]
\end{frame}


%%%%=================================%%%%
%% Contenido del curso
\begin{frame}{Contenido del curso}
\begin{enumerate}
  \item \textbf{Fundamentos de Business Analytics}
  \begin{itemize}
    \item Semana 01: Presentación del curso y metodología
    \item Semana 02: Introducción a LLMs
  \end{itemize} \vspace{0.15cm}

  \item \textbf{Introducción al lenguaje de programación R}
  \begin{itemize}
    \item Semana 03: Fundamentos de R
    \item Semana 04: Manipulación y visualización: \texttt{dplyr} + \texttt{ggplot2}
  \end{itemize} \vspace{0.15cm}

  \item \textbf{Proceso analítico y exploración de datos}
  \begin{itemize}
    \item Semana 05: Proceso analítico y tipos de analítica 
    \item Semana 06: EDA: fuentes, limpieza y exploración
  \end{itemize} \vspace{0.15cm}

  \item \textbf{Evaluación}
  \begin{itemize}
    \item Semana 07: Examen Parcial
  \end{itemize}
\end{enumerate}
\end{frame}


%%%%=================================%%%%
%% Contenido del curso (continuación)
\begin{frame}{Contenido del curso (continuación)}
\small
\begin{enumerate}
  \item \textbf{IA aplicada al análisis de datos}
  \begin{itemize}
    \item Semana 08: Claude Code / Cursor / VS Code
  \end{itemize}

  \item \textbf{Avances del proyecto}
  \begin{itemize}
    \item Semana 09: Presentación del avance del proyecto (EDA)
  \end{itemize}

  \item \textbf{Fundamentos de Machine Learning}
  \begin{itemize}
    \item Semana 10: Fundamentos de Machine Learning
  \end{itemize}  

  \item \textbf{Aprendizaje supervisado}
  \begin{itemize}
        \item Semana 11: Clasificación (pipeline, métricas y modelos base)
        \item Semana 12: Regresión (pipeline, métricas y modelos base)
  \end{itemize}  

  \item \textbf{Aprendizaje no supervisado}
  \begin{itemize}
    \item Semana 13
    : Clustering (fundamentos y métricas)
  \end{itemize}  

  \item \textbf{Evaluación}
  \begin{itemize}
        \item Semana 14: Examen Parcial
  \end{itemize}  

  \item \textbf{Simulacro}
  \begin{itemize}
        \item Semana 15: Simulacro de la presentación final
  \end{itemize}  
  
  \item \textbf{Cierre}
  \begin{itemize}
  \item Semana 16: Sustentación del proyecto final
  \end{itemize}
\end{enumerate}
\end{frame}


%%%%=================================%%%%
%% Agenda
\subsection{Dinámica de aprendizaje}
\begin{frame}[noframenumbering]{Agenda:}
\tableofcontents[currentsection,currentsubsection]
\end{frame}

%%%%=================================%%%%
%% Dinámica de aprendizaje
\begin{frame}{Dinámica de aprendizaje (cómo vamos a aprender)}
\small
\begin{itemize}
  \item \textbf{Estructura general:} Fundamentación teórica + Aplicación práctica. \vspace{0.15cm}

  \item \textbf{Antes de clase (pre-clase):} 
  \begin{itemize}
    \item Video corto / lectura guiada / podcast: Fundamentación teórica.
    \item Llegar con preguntas y dudas preparadas.
  \end{itemize} \vspace{0.15cm}

  \item \textbf{Inicio de cada clase:} quices conceptuales
  \begin{itemize}
    \item Se realizara cada semana un quiz a través de la plataforma Intu, sobre el material previo (podcast o video).
    \item Podran ser de selección multiple, verdadero/falso, ejercicios breves o interpretación de salidas
    \item La retroalimentación de cada quiz estara disponible en la plataforma para revisión del estudiante.
    \item \textbf{NO} hay recuperación por inasistencia; se elimina la peor calificación de quices una vez durante el semestre.
  \end{itemize} \vspace{0.15cm}
\end{itemize}
\end{frame}
%% Preguntar si poner como requisito los quices en los pcs de Icesi
%%%%=================================%%%%
%%%%=================================%%%%
%% Dinámica de aprendizaje
\begin{frame}{Dinámica de aprendizaje}
\small
\begin{itemize}
  \item \textbf{Durante la clase:}
  \begin{itemize}
    \item Taller guiado (\textit{hands-on}) en R / análisis de datos.
    \item Discusión: interpretación y decisión (qué significa el resultado y por qué importa).
  \end{itemize} \vspace{0.15cm}
  \item \textbf{Antes de finalizar la clase:}
  \begin{itemize}
    \item Actividad individual en clase.
    \item Repaso + práctica individual.
  \end{itemize}
\end{itemize}
\end{frame}

%%%%=================================%%%%
%% Estructura de Intu
\begin{frame}{Estructura de Intu}
\begin{figure}
    \centering
    \includegraphics[width=0.92\linewidth]{slides/week-01/figures/intu.png}
\end{figure}
\end{frame}

%%%%=================================%%%%
%% Agenda
\subsection{Evaluación y reglas de juego}
\begin{frame}[noframenumbering]{Agenda:}
\tableofcontents[currentsection,currentsubsection]
\end{frame}


%%%%=================================%%%%
%% Evaluación
\begin{frame}{Evaluación}
\small
\begin{itemize}
  \item \textbf{Talleres prácticos en clase (10\%)}
  \begin{itemize}
    \item Actividades \textbf{hands-on} desarrolladas durante la sesión.
    \item Evidencia evaluable: entrega de código, reporte corto, respuestas en plataforma o selección múltiple.
    \item Criterios típicos: ejecución correcta, manejo de datos, interpretación y trazabilidad de la evidencia.
  \end{itemize}

  \vspace{0.15cm}

  \item \textbf{Quices conceptuales (25\%)}
  \begin{itemize}
    \item Quices semanales durante el semestre; pueden ser anunciados o no anunciados.
    \item Formatos posibles: selección múltiple, V/F, respuesta corta o ejercicios breves.
    \item Evalúan precisión conceptual, razonamiento e interpretación de resultados.
  \end{itemize}

  \vspace{0.15cm}

  \item \textbf{Evaluaciones escritas integradoras: Parciales (40\%)}
  \begin{itemize}
    \item Parcial 1: \textbf{20\%} (contenidos de Unidades 1 a 3).
    \item Parcial 2: \textbf{20\%} (contenidos d Unidades 4 a 7).
  \end{itemize}

  \vspace{0.15cm}

  \item \textbf{Proyecto final aplicado progresivo / Unidad 8 (25\%)}
  \begin{itemize}
    \item Trabajo aplicado (en grupos de 3 estudiantes) con entregas progresivas.
    \item Se evalúa: selección del método, ejecución, interpretación y comunicación.
  \end{itemize}
\end{itemize}
\end{frame}


%%%%=================================%%%%
%% Reglas de juego
\begin{frame}{Reglas de juego (y supletorios)}
\small
\begin{itemize}
  \item \textbf{Reglas de juego (operativas)}
  \begin{itemize}
    \item Quices y talleres se realizan \textbf{únicamente en clase} (no hay reposición).
    \item Se elimina \textbf{una (1) nota} por criterio: la peor calificación de quices y la peor calificación de actividades en clase (según aplique).
    \item Asistencia mínima: \textbf{80\%}. Asistir a menos de este porcentaje es causal de no aprobación.
  \end{itemize} \vspace{0.2cm}

  \item \textbf{Política de supletorios (según reglamento ICESI)}
  \begin{itemize}
    \item La recuperación de notas aplica \textbf{únicamente para exámenes} (parcial,final y/o presentación final) cuando no se puedan presentar en la fecha programada por razones justificadas y debidamente soportadas.
    \item \textbf{No} se programan supletorios, reposiciones ni cambios de calificación para \textbf{quices, talleres o actividades en clase}, independientemente de la causa de inasistencia.
    \item La solicitud debe tramitarse con el director del programa (\textbf{no con el profesor}) dentro de los dos (2) días hábiles siguientes a la fecha del examen.
    \item Si la solicitud es aprobada: se genera autorización para pago en caja; luego se envía el comprobante a la auxiliar administrativa, quien informa programación y condiciones del supletorio.
  \end{itemize}
\end{itemize}
\end{frame}


%%%%=================================%%%%
%% Agenda
\subsection{Política de IA}
\begin{frame}[noframenumbering]{Agenda:}
\tableofcontents[currentsection,currentsubsection]
\end{frame}

%%%%=================================%%%%
%% Evidencia: efectos de IA en formación de habilidades
\begin{frame}{IA y formación de habilidades (Shen \& Tamkin, 2026)}
\small
\begin{itemize}
  \item \textbf{Pregunta del estudio:} ¿la asistencia con IA acelera el trabajo \textit{y} mejora (o deteriora) el aprendizaje cuando se aprende una habilidad nueva? \vspace{0.1cm}
  
  \item \textbf{Diseño:} experimento aleatorizado con desarrolladores aprendiendo una librería nueva (programación asíncrona en Python). \vspace{0.1cm}
  
  \item \textbf{Resultados principales:}
  \begin{itemize}
    \item Con IA, los participantes obtienen peor desempeño en evaluación de habilidades: cae el puntaje en 17\% (aprox. 2 notas) en \textit{comprensión conceptual, lectura de código y debugging}. 
    \item En promedio, no hay ganancia significativa de tiempo al completar las tareas.
    \item El mayor deterioro aparece en debugging (cuando se delega la solución de errores a la IA).
  \end{itemize} \vspace{0.1cm}
  
  \item \textbf{Mensaje para el curso:} productividad con IA no es atajo a competencia.
        La IA funciona mejor para aprender cuando el estudiante se mantiene cognitivamente activo.
        (p.ej., pedir \textit{explicaciones}, hacer \textit{preguntas conceptuales}, y verificar la lógica).
\end{itemize}

\vspace{0.15cm}
\footnotesize
\textit{Fuente: ``How AI Impacts Skill Formation'' (Shen \& Tamkin, 2026) \href{https://arxiv.org/abs/2601.20245}{https://arxiv.org/abs/2601.20245}}
\end{frame}


%%%%=================================%%%%
%% Política de IA
\begin{frame}{Política de IA (uso responsable y trazable)}
\small
\begin{itemize}
  \item \textbf{Principio:} la IA puede apoyar el aprendizaje, pero el estudiante debe entender y poder explicar lo entregado. \vspace{0.15cm}

  \item \textbf{Se permite (con transparencia)}
  \begin{itemize}
    \item Aclarar conceptos, generar ejemplos, revisar redacción.
    \item Ayuda para depurar código (\textit{debugging}) y entender errores.
    \item Sugerencias de estructura para reportes y presentaciones.
  \end{itemize} \vspace{0.15cm}

  \item \textbf{No se permite}
  \begin{itemize}
    \item Entregar como propio trabajo generado por IA sin citarlo/explicarlo.
    \item Usar IA para responder evaluaciones \textbf{cerradas} o \textbf{sin autorización} (quices, parciales y actividades en clase, salvo indicación explícita).
    \item Fabricar resultados, datos, tablas, citas o referencias (\textit{hallucinations}).
  \end{itemize} \vspace{0.15cm}

  \item \textbf{Regla práctica:} si no puedes explicar y reproducir tu resultado, no cuenta como trabajo válido.
\end{itemize}
\end{frame}


%%%%=================================%%%%
\subsection{Fechas clave}
\begin{frame}[noframenumbering]{Agenda:}
\tableofcontents[currentsection,currentsubsection]
\end{frame}


%%%%=================================%%%%
%% Agenda
\begin{frame}{Fechas clave}
\begin{itemize}
  \item Inicio de clases: \textbf{27 de julio de 2026}
  \item Terminación de clases: \textbf{14 de noviembre de 2026}
  \item Examen Parcial 1: \textbf{7 al 13 de septiembre de 2026}
  \item Presentación EDA: \textbf{21 al 27 de septiembre de 2026}
  \item Examen Parcial 2: \textbf{26 al 31 de octubre de 2026}
  \item Simulacro sustentación: \textbf{2 al 8 de noviembre de 2026}
  \item Sustentación trabajo final: \textbf{9 al 14 de noviembre de 2026}
  
\end{itemize}
\end{frame}

%% Charlar con Eduard cambiar el task de EDA por tiempo en clase para
%% Pensar en la pregunta y el enfoque

%%%%=================================%%%%
%% Agenda
\section{Proyecto aplicado del curso}
\begin{frame}[noframenumbering]{Agenda:}
\tableofcontents[currentsection]
\end{frame}

%%%%=================================%%%%
%% Trabajo final (Proyecto aplicado)
\begin{frame}{Trabajo final (Proyecto aplicado)}

\begin{enumerate}
\item \textbf{Entregable 1 (0\%):} \textit{Pregunta de negocio y roles}\vspace{0.1cm}
A partir de una base de datos y posibles temas sugeridos:
\begin{itemize}
\item Formular la \textit{pregunta de negocio} a resolver.
\item Proponer los \textit{roles} que se necesitarían en un caso real (p.\,ej., analista, ingeniero/a de datos, experto/a del negocio, etc.).
\end{itemize} \vspace{0.1cm}

\item \textbf{Entregable 2 (5\%):} \textit{Base depurada + descriptivas} \vspace{0.1cm}
\begin{itemize}
\item Limpieza y preparación de la base (variables, faltantes, formatos).
\item \textbf{Descriptivas} y exploración básica (tablas/gráficos clave).
\end{itemize} \vspace{0.1cm}

\item \textbf{Entrega final (20\%):} \textit{Sustentación + script} \vspace{0.1cm}
\begin{itemize}
\item \textit{Sustentación final (15\%)} en semana de exámenes.
\item \textit{Diapositivas (4\%)} con resultados y recomendaciones.
\item \textit{Script (1\%)} ordenado y comentado.
\end{itemize}
\end{enumerate}
\end{frame}


%%%%=================================%%%%
%% Inscripción de grupos (QR)
\begin{frame}{Inscripción de grupos (Trabajo final)}

\begin{itemize}
\item Escanea el \textbf{QR} para inscribir tu grupo.
\item Un registro por grupo.
\item Completa: \textit{integrantes y código de estudiante}.
\item Fecha límite: viernes 28 de agosto a las 23:59
\end{itemize}
\centering
\vspace{0.2cm}
\includegraphics[width=0.45\linewidth]{slides/week-01/figures/proyect_final.png}
\end{frame}

%%%%=================================%%%%
%% Cierre
\begin{frame}
\vspace{2.0cm}
\begin{center}
    \LARGE{Gracias}
\end{center}

\end{frame}

%%%%%%%%%%%%%%%%%%%%%%%%%%%%%%%%%%%%
\end{document} 